\begin{frame}
  \frametitle{Morris-Pratt algorithm}

In case of mismatch, the na\"{\i}ve algorithm resumes comparing the
first letters of~\(p\) without using the information of the partial
success, to wit, we know \(\ind{p}{0,i-1} = \ind{t}{j-i,j-1}\) and
\(\ind{p}{i} \neq \ind{t}{j}\).
\[
\begin{array}{@{}rcc@{\,}c@{}ccr@{}r@{\,}c@{}ccc@{}ccc@{}}
  & & & \scriptstyle j-i
  & & &
  & & \multicolumn{2}{@{}l}{\scriptstyle j-i+k}
  & &
  & \multicolumn{1}{r@{}}{\scriptstyle j-1}
  & \multicolumn{1}{c}{\scriptstyle j}\\
\cline{2-15}
    \multicolumn{1}{r|}{t:}
  & \multicolumn{2}{c}{\phantom{H}\ldots\phantom{H}}
  & \multicolumn{1}{|c|}{a}
  & \multicolumn{1}{c|}{b}
  & \multicolumn{1}{c|}{c}
  & \multicolumn{2}{c|}{\ldots}
  &
  &
  &
  & \phantom{HHHH}
  &
  & \multicolumn{1}{|>{\columncolor{mygrey}}c|}{\textcolor{white}{\texttt{b}}}
  & \phantom{HHHH}\\
\cline{2-15}
  &
  &
  & \multicolumn{1}{:c}{\scriptstyle 0}
  &
  &
  &
  &
  & \multicolumn{1}{:l@{\,}}{\scriptstyle i-k}
  &
  &
  &
  & \multicolumn{1}{c@{}:}{\scriptstyle i-1}
  & \multicolumn{1}{c}{\scriptstyle i}\\
\cline{4-15}
    \multicolumn{1}{r}{p:}
  &
  &
  & \multicolumn{1}{|c|}{a}
  & \multicolumn{1}{c|}{b}
  & \multicolumn{1}{c|}{c}
  & \multicolumn{2}{c}{\ldots}
  &
  &
  &
  &
  &
  & \multicolumn{1}{|>{\columncolor{mygrey}}c|}{\textcolor{white}{\texttt{a}}}
  & \\
\cline{4-15}
  &&&&&&
  &
  & \multicolumn{1}{:c}{\scriptstyle 0}
  &
  &
  &
  & \multicolumn{1}{c:}{}\\
\cline{9-15}
  &&&&&
  & \text{shifted}
  & \multicolumn{1}{c}{p:}
  & \multicolumn{1}{|c|}{a}
  & \multicolumn{1}{c|}{b}
  & \multicolumn{1}{c|}{c}
  & \multicolumn{2}{c}{\ldots}\\
\cline{9-15}
\end{array}
\]

\bigskip

If a new search starts in \(t\) at position \(j-i+k\), with \(1
\leqslant k \leqslant i\), it compares a prefix of~\(p\) to \(t[j-i+k
  \dots]\). Since \(t[j-i+k \dots j-1] = p[i-k,i-1]\), it compares a
prefix of~\(p\) to a suffix of \(p[\dots i-1]\), that is, to a part
of itself!

\end{frame}

\begin{frame}
  \frametitle{Morris-Pratt algorithm}

The important point here is that these comparisons are
\emph{independent} of the text~\(t\), since they apply to part of the
pattern~\(p\) itself.

\bigskip

Contrary to na\"{\i}ve factoring, letters in the text are compared in
a strictly increasing order (never having to backtrack).

\begin{figure}
\centering
\includegraphics[bb=74 621 367 717]{mp}
\end{figure}

where \(\Border{}{p}\) is the \textbf{maximum border} of~\(p\): the
longest prefix of \(p\) that is also a suffix of \(p[\dots i-1]\).

\end{frame}


\begin{frame}
\frametitle{Morris-Pratt algorithm}

The strategy devised here is to \textbf{compute the maximum borders of
  all prefixes of~\(p\) \emph{before} we start the search} in the
text, and then use this information to avoid repeating the same
comparisons at different positions in the text.

\bigskip

This \textbf{preprocessing} on~\(p\) will indeed allow to recognise
the sole configurations where the search may succeed.

\bigskip

In other words, we need to understand borders.

\end{frame}

\begin{frame}
  \frametitle{Morris-Pratt algorithm}

But before that, let us consider an example where, in the end,
\(p\)~is not found to be a factor of~\(t\). As usual, letters on a
grey background correspond to mismatches. It is clear that
\(\Border{}{a} = \varepsilon\), for all letters~\(a\).
\begin{figure}
\centering
\includegraphics{mp_ex}
\end{figure}

(The text is not shown in its entirety.)

\end{frame}

\begin{frame}
  \frametitle{Morris-Pratt algorithm}

The \textbf{border} of a non\hyp{}empty word~\(y\) is a prefix
of~\(y\) which is also a suffix.

\bigskip

To understand the case of the \textbf{empty border}, consider
\(\Border{}{\word{abac}} = \varepsilon\), simply because we have
\(\word{abac} =
\underline{\varepsilon}\word{abac}\underline{\varepsilon}\).

\bigskip

For example, the word `\word{abacaba}' has three borders:
\(\varepsilon\), `\word{a}'~and~`\word{aba}'. The last one is the
\textbf{maximum border} and we write
\(\Border{}{\underline{\word{aba}}\word{c}\underline{\word{aba}}} =
\word{aba}\).

\bigskip

Note that \textbf{borders can overlap}; for example,
\(\Border{}{\underline{\word{aaa}}\word{a}} =
\Border{}{\word{a}\underline{\word{aaa}}} = \word{aaa}\).

\end{frame}

\begin{frame}
  \frametitle{Morris-Pratt algorithm}

As often when working at the level of specifications, we seek an
inductive definition for the function at hand, in this instance, the
maximum border of a word.

\bigskip

We then have the choice of finding either \(\Border{}{ay}\) or
\(\Border{}{ya}\), with \(y \neq \varepsilon\), assuming that we know
\(\Border{}{y}\).

\bigskip

Since we are interested in knowing the maximum borders of all the
\emph{prefixes} of a given pattern, the \(ya\) is more suitable (\(y
\pref ya\)).

\bigskip

The idea is to recursively consider \(\Border{}{y} \cdot a\):

\begin{itemize}

  \item if it is a prefix of~\(y\), then \(\Border{}{ya} =
    \Border{}{y} \cdot a\);

  \item otherwise, we consider the border of \(\Border{}{y} \cdot a\).

\end{itemize}

\end{frame}


\begin{frame}
  \frametitle{Morris-Pratt algorithm}

Consider the following examples where \(y\)~and~\(\Border{}{y}\) are
given:
\begin{align*}
  y             &= \word{abaabb},
& \Border{}{y}  &= \varepsilon,
& \Border{}{y \cdot \word{b}}
                &= \Border{}{\Border{}{y} \cdot \word{b}}
                 = \Border{}{\word{b}} = \varepsilon;\\
  y             &= \word{baaaba},
& \Border{}{y}  &= \word{ba},
& \Border{}{y \cdot \word{a}}
                &= \Border{}{y} \cdot \word{a}
                 = \word{baa};\\
  y             &= \word{abbbab},
& \Border{}{y}  &= \word{ab},
& \Border{}{y \cdot \word{a}}
                &= \Border{}{\Border{}{y} \cdot \word{a}}
                 = \Border{2}{y} \cdot \word{a}
                 = \word{a}.
\end{align*}

For example,
\(\Border{}{y \cdot \word{a}}
   = \Border{}{\Border{}{y} \cdot \word{a}}
   = \Border{}{\Border{2}{y} \cdot \word{a}}
   = \Border{3}{y} \cdot \word{a}\):
\begin{figure}
\centering
\includegraphics{max_B}
\end{figure}

\end{frame}

\begin{frame}
  \frametitle{Morris-Pratt algorithm}

Formally, for all words~\({y \neq \varepsilon}\) and any
letter~\(a\),
\begin{equation*}
  \Border{}{a}         := \varepsilon;\qquad
  \Border{}{y \cdot a} := \left\{
    \begin{aligned}
      & \Border{}{y} \cdot a,
      && \text{if \(\Border{}{y} \cdot a \prefeq y\)};\\
      & \Border{}{\Border{}{y} \cdot a},
      && \text{otherwise.}
    \end{aligned}
  \right.
\end{equation*}

Unfolding the equations yields:
\begin{align*}
   \Border{}{ya}
&= \Border{}{\Border{}{y} \cdot a},
& \Border{}{y} \cdot a &\nprefeq y;\\
   \Border{}{\Border{}{y} \cdot a}
&= \Border{}{\Border{2}{y} \cdot a},
&  \Border{2}{y} \cdot a &\nprefeq \Border{}{y};\\
&\;\;\smash{\vdots} & &\;\;\smash{\vdots}\\
   \smash{\Border{}{\Border{p-1}{y} \cdot a}}
&= \smash{\Border{}{\Border{p}{y} \cdot a}},
& \smash{\Border{p}{y} \cdot a} &\nprefeq \smash{\Border{p-1}{y}};
\end{align*}
and \(\varepsilon \not\in \{y,\Border{}{y}, \dots,
\Border{p-1}{y}\}\).

\end{frame}


\begin{frame}
  \frametitle{Morris-Pratt algorithm}

Therefore, an equivalent definition is:
\begin{equation*}
\Border{}{ya} =
\left\{
  \begin{aligned}
   & \Border{q}{y} \cdot a,
   && \text{if \(\Border{q}{y} \cdot a \prefeq y\)};\\
   & \varepsilon,
   && \text{otherwise;}
  \end{aligned}
\right.
\end{equation*}
with the additional constraint that \textbf{\(q\)~must be as small as
  possible} (so we define the maximum border).

\bigskip

This form of the definition of~\(\BorderName\) is simpler because it
does not contain an embedded call.

\end{frame}


\begin{frame}
  \frametitle{Morris-Pratt algorithm}

We do not actually need to compute the maximum borders of all the
prefixes of the pattern. Let us have a look back at
\begin{figure}
\centering
\includegraphics[bb=74 621 367 717]{mp}
\end{figure}

What we need is the offset corresponding to the \textbf{length of the
  maximum border} being considered (\(k\), in the figure).

\end{frame}


\begin{frame}
  \frametitle{Morris-Pratt algorithm}

Let us define a function \(\MPfailureName_{x}\) on all the prefixes of
a word~\(x\) as follows:
\begin{equation}
  \MPfailure{x}{}{\len{y}}
:= \len{\Border{}{y}},\,\; \text{for
    all \(x\) and \(y \neq \varepsilon\) such
    that \(y \prefeq x\)}.\label{eq:MP_failure}
\end{equation}
For reasons which will be clear soon, this function is called the
\textbf{failure function} of~\(x\). An equivalent definition is
\begin{equation*}
  \MPfailure{x}{}{i}
= \len{\Border{}{\ind{x}{\dots i-1}}}, \,\;
\text{for all \(x\) and  \(i\) such that
\(0 < i \leqslant \len{x}\)}.
\end{equation*}

Let's take the equations deining of the maximum borders and take the
lengths on both sides:

\begin{equation*}
\len{\Border{}{ya}} =
\left\{
  \begin{aligned}
   & \len{\Border{q}{y} \cdot a} = 1 + \len{\Border{q}{y}},
   && \text{if \(\Border{q}{y} \cdot a \prefeq y\)};\\
   & \len{\varepsilon} = 0,
   && \text{otherwise.}
  \end{aligned}
\right.
\end{equation*}

\end{frame}


\begin{frame}
  \frametitle{Morris-Pratt algorithm}

If~\({ya \prefeq x}\), then \(\len{\Border{}{ya}} =
\MPfailure{x}{}{\len{ya}} = \MPfailure{x}{}{\len{y}+1}\). Let \(i :=
\len{y} > 0\).
\begin{equation*}
\MPfailure{x}{}{i+1} =
\left\{
  \begin{aligned}
   & 1 + \len{\Border{q}{y}},
   && \text{if \(\Border{q}{y} \cdot a \prefeq y\)};\\
   & 0,
   && \text{otherwise.}
  \end{aligned}
\right.
\end{equation*}

Let's set aside for a moment \(\Border{q}{y} \cdot a \prefeq y\) and
return to the definition of~\(\MPfailureName\), from which we deduce:
\begin{equation*}
\MPfailure{x}{q}{\len{y}} = \len{\Border{q}{y}},
\,\; \text{with \(y \prefeq x\)},
\end{equation*}
which we can prove by complete induction on~\(q\). Let us call this
property \(\pred{P}{q}\). Trivially, we have \(\pred{P}{0}\). Let us
suppose \(\pred{P}{n}\) for all \(n \leqslant q\): this is the
induction hypothesis. We suppose \(y \prefeq x\) and prove now
\(\pred{P}{q+1}\):
\begin{equation*}
  \MPfailure{x}{q+1}{\len{y}}
= \MPfailure{x}{q}{\MPfailure{x}{}{\len{y}}}
= \MPfailure{x}{q}{\len{\Border{}{y}}}
\doteq \len{\Border{q}{\Border{}{y}}}
= \len{\Border{q+1}{y}},
\end{equation*}
where \((\smash{\doteq})\)~is a valid application of the induction
hypothesis because \(\Border{}{y} \pref y \prefeq x\). This proves
\(\pred{P}{q+1}\) and the induction principle entails that
\(\pred{P}{n}\) holds for all \(n \geqslant 0\).

\end{frame}


\begin{frame}
  \frametitle{Morris-Pratt algorithm}

  So
  \begin{equation*}
\MPfailure{x}{}{i+1} =
\left\{
  \begin{aligned}
   & 1 + \MPfailure{x}{q}{i},
   && \text{if \(\Border{q}{y} \cdot a \prefeq y\)};\\
   & 0,
   && \text{otherwise.}
  \end{aligned}
\right.
\end{equation*}
There is one part of the definition that we did not use:
\({\Border{}{a} := \varepsilon}\). It implies \(\MPfailure{x}{}{1} =
\MPfailure{x}{}{\len{a}} = \len{\Border{}{a}} = \len{\varepsilon} =
0\) and, since the definition of~\(\MPfailureName\) implies
\(\MPfailure{x}{}{1} = 1 + \MPfailure{x}{}{0}\), so
\(\MPfailure{x}{}{0} = -1\).

\bigskip

We can now return to the last missing bit:
\begin{equation*}
\Border{q}{y} \cdot a \prefeq y
\,\; \text{and} \,\; ya \prefeq x \,\;\text{and}\,\; \len{y} = i.
\end{equation*}
It is equivalent to the following:
\begin{equation*}
  \smash{\Ind{y}{\len{\Border{q}{y}}}} = a \Leftrightarrow
\smash{\Ind{y}{\MPfailure{x}{q}{i}}} = a \Leftrightarrow
\smash{\Ind{x}{\MPfailure{x}{q}{i}}} = \smash{\Ind{x}{\len{y}}}
\Leftrightarrow \smash{\Ind{x}{\MPfailure{x}{q}{i}}} = \ind{x}{i}.
\end{equation*}

\end{frame}


\begin{frame}
  \frametitle{Morris-Pratt algorithm}

Gathering everything:
\begin{equation*}
 \MPfailure{x}{}{0}   = -1\quad\text{and}\quad
 \MPfailure{x}{}{i+1} =
   \left\{
     \begin{aligned}
       & 1 + \MPfailure{x}{q}{i},
       && \text{if \(\Ind{x}{\MPfailure{x}{q}{i}} = \ind{x}{i}\)};\\
       & 0,
       && \text{otherwise;}
     \end{aligned}
   \right.
\end{equation*}
where \(q\)~is the smallest nonzero natural satisfying the
condition.

\bigskip

This can be further simplified into
\begin{equation*}
\MPfailure{x}{}{0} = -1
\quad \text{and} \quad
\MPfailure{x}{}{i+1} = 1 + \MPfailure{x}{q}{i}, \,\;\text{with}\,\; i \geqslant 0,
\end{equation*}
and \(q>0\)~is the smallest natural such that
\(\MPfailure{x}{q}{i} = -1\) or \(\smash{\Ind{x}{\MPfailure{x}{q}{i}}}
= \ind{x}{i}\).

\end{frame}
